\begin{Definition}
\label{def:symmetric-operator}
对于从 $H$ 到 $H$ 且定义域 $\mathcal{D}(A)$ 稠密的算子 $(\mathcal{D}(A),A)$, 若 $(\mathcal{D}(A^*),A^*)$ 是 $(\mathcal{D}(A),A)$ 的扩张, 或者等价地, 若
\begin{equation}\label{eq:symmetric-operator-pairing-identity}
(Au,v)=(u,Av)\quad\text{for every }u,v\in\mathcal{D}(A).
\end{equation}
则称该算子是\textbf{对称的}.
\end{Definition}

\begin{Definition}
\label{def:self-adjoint-operator}
若算子满足
\begin{equation}\label{eq:self-adjoint-operator-domain-action-equality}
(\mathcal{D}(A),A)=(\mathcal{D}(A^*),A^*).
\end{equation}
则称它是\textbf{自伴随的}.
\end{Definition}

以下引理是定义的直接推论.

\begin{Lemma}
\label{lem:self-adjoint-symmetric-domain-equality}
算子 $(\mathcal{D}(A),A)$ 是自伴随的, 当且仅当它是对称的且 $\mathcal{D}(A)=\mathcal{D}(A^*)$.
\end{Lemma}

对于有界算子, 对称性与自伴随性是一回事.

\begin{Theorem}
\label{thm:bounded-symmetric-operator-self-adjoint}
算子 $A\in\mathcal{L}(H)$ 是自伴随的, 当且仅当它是对称的.
\end{Theorem}

\begingroup
\postdisplaypenalty=10000
\renewcommand{\endProof}{\par\nobreak\hfill$\blacksquare$\par\vspace{0.1cm}}
\begin{Proof}
由定义, 自伴随性蕴含对称性. 此外, 若 $\mathcal{D}(A)=H$ 且 $A$ 是对称的, 则由 \eqref{eq:symmetric-operator-pairing-identity} 得 $\mathcal{D}(A^*)=\mathcal{D}(A)=H$.
\end{Proof}
\endgroup

正如在 Problem~\ref{ex:second-derivative-symmetric-not-self-adjoint} 中将看到的, 存在并非自伴随的无界对称算子.

\Needspace{4\baselineskip}
\begin{Example}
\label{ex:differentiation-hilbert-adjoint-complete-calculation}
现在给出一个计算无界算子伴随的例子. 设 $H=L^2(0,1)$, 并令\[\mathcal{D}(D)=\{u\in H^1(0,1)\mid u(0)=0\}.\]这里的边界条件在迹的意义下理解. 对所有 $u\in\mathcal{D}(D)$, 用下式定义从 $H$ 到 $H$ 的微分算子 $(\mathcal{D}(D),D)$:\[D(u):=u'\]先作形式计算, 对伴随的具体形式给出一个 ``猜测'', 然后严格证明这一猜测正确.

形式计算和严格证明均基于恒等式
\begin{equation}\label{eq:differentiation-adjoint-defining-pairing}
(v,Du)=(D^*v,u)\quad\text{for all }u\in\mathcal{D}(D),\ v\in\mathcal{D}(D^*).
\end{equation}
在形式计算中, 从 \eqref{eq:differentiation-adjoint-defining-pairing} 的左边出发作分部积分 (尽管对于 $v$ 目前除了它属于 $L^2(0,1)$ 外尚无任何了解). 得到\[\begin{aligned}
(v,Du)&=\int_0^1u'(x)v(x)\,\mathrm{d}x\\
&=-\int_0^1u(x)v'(x)\,\mathrm{d}x+u(1)v(1)-u(0)v(0)\\
&=(u,-v')+u(1)v(1).
\end{aligned}\]最后一个等式使用了 $u(0)=0$ 这一事实. 考察 \eqref{eq:differentiation-adjoint-defining-pairing} (并要求形式计算有意义), 可猜测伴随由 $(\mathcal{D}(B),B)$ 给出, 其中
\begin{equation}\label{eq:differentiation-adjoint-candidate-domain-action}
\mathcal{D}(B):=\{v\in H^1(0,1)\mid v(1)=0\},\quad B(v)=-v'.
\end{equation}
这一猜测是正确的, 但目前只证明了伴随是算子 $(\mathcal{D}(B),B)$ 的\emph{扩张}. 为证明 $(\mathcal{D}(D^*),D^*)=(\mathcal{D}(B),B)$, 令 $D^*v:=f$ (除了属于 $L^2(0,1)$ 外不作其他假设), 并定义
\begin{equation}\label{eq:differentiation-adjoint-antiderivative}
F(x)=\int_x^1f(s)\,\mathrm{d}s.
\end{equation}
当然, 由 \eqref{eq:differentiation-adjoint-candidate-domain-action} 可知 $F\in\mathcal{D}(B)$ 且 $B(F)=f$. 现在从 \eqref{eq:differentiation-adjoint-defining-pairing} 的\emph{右边}出发, 对每个 $v\in\mathcal{D}(D^*)$ 和 $u\in\mathcal{D}(D)$ 得到\[\begin{aligned}
(D^*v,u)&=(f,u)\\
&=\int_0^1u(x)f(x)\,\mathrm{d}x\\
&=-\int_0^1u(x)F'(x)\,\mathrm{d}x\\
&=\int_0^1u'(x)F(x)\,\mathrm{d}x\\
&=(F,Du)\\
&=(v,Du).
\end{aligned}\]由于这对所有 $u\in\mathcal{D}(D)$ 成立, 必有 $F-v\in\mathcal{R}(D)^\perp$. 然而, $\mathcal{R}(D)=L^2(0,1)$, 因而 $F=v$. 断言得证.
\end{Example}

\Needspace{4\baselineskip}
\begin{Example}
\label{ex:general-differential-hilbert-formal-adjoint}
对于更一般的偏微分算子, 例如 Example~\ref{ex:differential-operator-boundary-domains} 中定义的算子, 由于目前还无法计算算子的闭包, 伴随的计算变得复杂. 然而, 无需先使算子闭合, 就能证明 Hilbert 伴随是 Definition~\ref{lem:distribution-convolution-continuity} 中定义的形式伴随的\emph{扩张}.
\end{Example}

正如下文将看到的, 自伴随算子具有许多在应用中有用的性质. 然而, 正如上述例子所表明的, 判断无界算子是否自伴随往往很困难. (必须先使算子闭合, 而且如前所见, 求闭包的定义域往往并非易事.) 幸运的是, 可以定义一种称为\emph{本质自伴随性}的相关性质, 它有相对容易的判别方法.
